\section{Discussion}
\label{sec:discussion}

\subsection{Key Findings}

\textbf{The bias is large, structural, and persistent.} Cram{\'e}r's $V = 0.40$--$0.57$ places the language $\times$ retention association firmly in the ``large'' effect range. The consistent \textit{es} $>$ \textit{fr} $>$ \textit{it} $>$ \textit{en} $>$ \textit{de} ordering across all five threshold levels is not compatible with a sampling artifact; it tracks the Romance--Germanic language family boundary precisely. At $k$=30, a practitioner would retain 6.3$\times$ more Spanish documents than German documents from the same raw corpus. The training data this produces does not reflect relative document quality; it reflects a calibration artifact.

\textbf{Prior estimates underestimate the bias.} Our finding that $V = 0.40$--$0.57$ exceeds literature-derived estimates by 25--40\% has direct consequences for correction study design. We recommend using $V = 0.40$--$0.57$ as the calibrated baseline when designing correction experiments on the RedPajama-V2 \texttt{ccnet\_perplexity} signal.

\textbf{Connection to prior work.} Our measurement extends \citet{Caswell2021Quality}'s qualitative documentation to a specific quantitative baseline. It confirms CCNet's original per-language design intent \cite{Wenzek2019CCNet} by precisely quantifying what happens when that intent is not followed. It provides the $V$ baseline that \citet{Turki2026CrossLingual}'s endorsement of percentile-based per-language thresholding implicitly requires for evaluation.

\subsection{Limitations}

\textbf{L1: Existence-Only Scope.} This paper tests only the existence of the disparity under global thresholding. Whether per-language calibration reduces $\Delta V \geq 0.10$ for $\geq 3/5$ $k$ values remains to be tested. The characterization contribution is complete and independently publishable.

\textbf{L2: Sample Scope.} Results are based on 208,262 documents ($\sim$0.0002\% of full 113.3B document corpus). The phenomenon is structurally driven and confirmed across three independent runs; extrapolation to the full corpus requires distributed-scale replication.

\textbf{L3: Document-Length Confound.} Longer documents tend to have lower perplexity. Length stratification was planned but not implemented. The 72.7pp gap substantially exceeds what a plausible length confound can explain; length stratification is a recommended robustness check.

\textbf{L4: Five European Languages.} The five-language sample covers only Germanic and Romance families. Findings may not generalize to low-resource or non-Indo-European languages.

\subsection{Broader Impact}

Any model trained on a corpus filtered using global percentile thresholds on \texttt{ccnet\_perplexity} will have structurally imbalanced language coverage. Our work enables practitioners to (a) audit existing multilingual training datasets for this bias source, (b) design properly powered correction studies, and (c) adopt per-language calibration as a default. We recommend evaluating correction strategies against the $V$ baseline established here before deploying in production pipelines.
